\begin{lemma}
\label{lemma-cokernel-flat}
Let $R$ be a Noetherian domain. Let $R \to A \to B$ be finite type ring maps.
Let $M$ be a finite $A$-module and let $N$ a finite $B$-module.
Let $M \to N$ be an $A$-linear map. There exists an nonzero $f \in R$
such that the cokernel of $M_f \to N_f$ is a flat $R_f$-module.
\end{lemma}

\begin{proof}
By replacing $M$ by the image of $M \to N$, we may assume $M \subset N$.
Choose a filtration $0 = N_0 \subset N_1 \subset \ldots \subset N_t = N$
such that $N_i/N_{i - 1} = B/\mathfrak q_i$ for some prime ideal
$\mathfrak q_i \subset B$, see
Algebra, Lemma \ref{algebra-lemma-filter-Noetherian-module}.
Set $M_i = M \cap N_i$. Then $Q = N/M$ has a filtration by the submodules
$Q_i = N_i/M_i$. It suffices to prove $Q_i/Q_{i - 1}$ becomes flat
after localizing at a nonzero element of $f$ (since extensions of
flat modules are flat by Algebra, Lemma \ref{algebra-lemma-flat-ses}).
Since $Q_i/Q_{i - 1}$ is isomorphic to the  cokernel
of the map $M_i/M_{i - 1} \to N_i/N_{i - 1}$, we reduce to the
case discussed in the next paragraph.

\medskip\noindent
Assume $B$ is a domain and $M \subset N = B$. After replacing $A$ by the
image of $A$ in $B$ we may assume $A \subset B$. By generic flatness,
we may assume $A$ and $B$ are flat over $R$
(Algebra, Lemma \ref{algebra-lemma-generic-flatness-Noetherian}).
It now suffices to show $M \to B$ becomes $R$-universally
injective after replacing $R$ by a principal localization
(Algebra, Lemma \ref{algebra-lemma-ui-flat-domain}).
By generic freeness, we can find a nonzero $g \in A$ such that
$B_g$ is a free $A_g$-module
(Algebra, Lemma \ref{algebra-lemma-generic-flatness-Noetherian}).
Thus we may choose a direct summand $M' \subset B_g$ as an $A_g$-module,
which is finite free as an $A_g$-module, and
such that $M \to B \to B_g$ factors through $M'$.
Clearly, it suffices to show that $M \to M'$
becomes $R$-universally injective after replacing
$R$ by a principal localization.

\medskip\noindent
Say $M' = A_g^{\oplus n}$. Since $M \subset M'$ is a finite $A$-module,
we see that $M$ is contained in $(1/g^m)A^{\oplus n}$ for some $m \geq 0$.
After changing our basis for $M'$ we may assume $M \subset A^{\oplus n}$.
Then it suffices to show that $A^{\oplus n}/M$ and $A_g/A$ become
$R$-flat after replacing $R$ by a principal localization. Namely, then
$M \to A^{\oplus n}$ and $A^{\oplus n} \to A_g^{\oplus n}$ are
universally injective by Algebra, Lemma \ref{algebra-lemma-flat-tor-zero}
and consequently so is the composition $M \to M' = A_g^{\oplus n}$.

\medskip\noindent
By generic flatness (see reference above), we may assume the
module $A^{\oplus n}/M$ is $R$-flat. For the quotient $A_g/A$
we use the fact that
$$
A_g/A = \colim (1/g^m)A/A \cong \colim A/g^mA
$$
and the module $A/g^mA$ has a filtration of length $m$ whose
successive quotients are isomorphic to $A/gA$. Again by generic
flatness we may assume $A/gA$ is $R$-flat and hence each $A/g^mA$
is $R$-flat, and hence so is $A_g/A$.
\end{proof}